\scanpage{473}
\mspage{539}

\begin{equation*}
\tag{34}
\left\{
\begin{gathered}
\diagram{\begin{tikzcd}[ampersand replacement=\&]
S\mathcal M_{\rho,\sigma}\arrow[r,"\sim"]\&
\mathcal M_{\rho,\sigma}\mathbin{\underset{\mathbb{G}_m}{\times}}B'_0
\end{tikzcd}}\\
\text{where}\quad\left\{
\begin{gathered}
\diagram{\begin{tikzcd}[ampersand replacement=\&]
\mathcal M_{\rho,\sigma}\arrow[r,"\chi"]\&\mathbb{G}_m\\
B'_0\arrow[r,"{\delta_{B'_0}=\delta_G|B'_0}"]\&\mathbb{G}_m
\end{tikzcd}}
\end{gathered}\right.
\end{gathered}\right.\quad,\quad\text{Whew !}
\end{equation*}

This formula will serve us as the \uline{definition} of
$S\mathcal M_{\rho,\sigma}$ from now on -- it defines at once
a group law on $S\mathcal M_{\rho,\sigma}$, and we find on
$S\mathcal M_{\rho,\sigma}$ \uline{two} extension structures

\begin{equation*}
\tag{35}
\diagram{\begin{tikzcd}[ampersand replacement=\&,column sep=small]
1\arrow[r]\&L_0\arrow[r]\arrow[d,phantom,"\Vert"]\&
S\mathcal M_{\rho,\sigma}\arrow[r]\&\mathcal M_{\rho,\sigma}\arrow[r]\&1\,,\\
\&\operatorname{Ker}(B'_0\xrightarrow{\delta_{B'_0}}\mathbb{G}_m)\&
\widetilde x=(x,\unsure{\underline u|y})\arrow[r,mapsto]\&x
\end{tikzcd}}
\end{equation*}

\begin{equation*}
\tag{36}
\diagram{\begin{tikzcd}[ampersand replacement=\&,column sep=small]
1\arrow[r]\&\mathcal M^+_{\rho,\sigma}\arrow[r]\arrow[d,phantom,"\Vert"]\&
S\mathcal M_{\rho,\sigma}\arrow[r]\&B'_0\arrow[r]\&1\\
\&\operatorname{Ker}(\mathcal M_{\rho,\sigma}\xrightarrow{\chi}\mathbb{G}_m)\&
\begin{gathered}\widetilde x\\\Vert\\(x,\underline u)\end{gathered}
\arrow[r,mapsto]\&\underline u
\end{tikzcd}}
\end{equation*}

The subgroup $S\mathcal M_{\rho,\sigma}[0]$ of
$S\mathcal M_{\rho,\sigma}$ corresponds to the condition $f=1$ i.e.
$\begin{gathered}\underline U\\\Vert\\\Theta_G(x)\end{gathered}=\underline u$
(which already implies
$\begin{gathered}\det\underline U\\\Vert\\\chi(x)\end{gathered}=\det\underline u$
-- it is also the subgroup of the $(x,u)$ in
$\mathcal M_{\rho,\sigma}\times B'_0$ such that we have
$u=\Theta_G(x)$. Thus $S\mathcal M_{\rho,\sigma}[0]$ is the
subgroup of $S\mathcal M_{\rho,\sigma}$ which is the
\inserted{isomorphic} image of $\mathcal M_{\rho,\sigma}[0]$
\deleted{\unsure{under}}
\inserted{(cf. (\unsure{14|19}) page 536')}
under $x\longmapsto(x,\Theta_G(x))$ :

\deleted{\diagram{\begin{tikzcd}[ampersand replacement=\&]
x\arrow[r,mapsto]\&\unsure{(}
\end{tikzcd}}}

\begin{equation*}
\tag{37}
S\mathcal M_{\rho,\sigma}[0]=\operatorname{Im}\left(
\diagram{\begin{tikzcd}[ampersand replacement=\&]
\mathcal M_{\rho,\sigma}[0]\arrow[r,hook]\&S\mathcal M_{\rho,\sigma}\\
x\arrow[r,mapsto]\&(x,\Theta_G(x))
\end{tikzcd}}\right)
\end{equation*}

\[
(\text{where}\quad\mathcal M_{\rho,\sigma}[0]
=\Theta_G^{-1}(B'_0)\quad)
\]

\mspage{472}

\scanpage{474}

Giving a splitting of (35) amounts to a splitting of
\[
\diagram{\begin{tikzcd}[ampersand replacement=\&]
1\arrow[r]\&L_0\arrow[r]\&B'_0\arrow[r]\&\mathbb{G}_m\arrow[r]\&1\,,
\end{tikzcd}}
\]
now there is one ready-made via the subgroup $T_0$. So we deduce from this
an embedding of $\mathcal M_{\rho,\sigma}$ into
$S\mathcal M_{\rho,\sigma}$, which thereby does indeed become a
$1/2$-direct product.

\begin{equation*}
\tag{38}
S\mathcal M_{\rho,\sigma}\simeq\mathcal M_{\rho,\sigma}\cdot L_0
\end{equation*}

$\mathcal M_{\rho,\sigma}$ \unsure{operating} moreover on $L_0$ via
the character $\chi:\mathcal M_{\rho,\sigma}\longrightarrow \mathbb{G}_m$ --
but this description now seems to me much less
``telling'' than (34), for it does not bring out the canonical
homomorphism
\[
\diagram{\begin{tikzcd}[ampersand replacement=\&]
S\mathcal M_{\rho,\sigma}\arrow[r]\&B'_0
\end{tikzcd}}
\]
which is nevertheless the essential thing for the introduction of
$S\mathcal M_{\rho,\sigma}$. Once this is in place, we find
(``for the record'') the

\uline{Theorem}\quad We have a canonical isomorphism

\begin{equation*}
\tag{39}
\diagram{\begin{tikzcd}[ampersand replacement=\&]
S\mathcal M_{\rho,\sigma}\arrow[r]\&
S\deleted{\unsure{\mathcal M}}\inserted{\Sigma}_{\rho,\sigma}^{*}
\end{tikzcd}}
\quad
\begin{gathered}
\text{group}\\
\text{scheme of the}\\
\text{systems }(\alpha,\beta,f)\\
\text{satisfying }(17)\\
\text{and }(31)
\end{gathered}
\end{equation*}

\deleted{\diagram{\begin{tikzcd}[ampersand replacement=\&]
(x,\underline u)\arrow[r,mapsto]\&
(\alpha=\underline u\unsure{\underline a|\Theta_\rho}^{-1},
\beta=\underline u\underline b^{-1},f=\underline u\underline U^{-1})
\end{tikzcd}}}

\deleted{where $\underline a=\Theta_\rho(x)$,
$\underline b=\Theta_\sigma(x)$, $\underline U=\Theta_G(x)$}

\mspage{473}

\scanpage{475}
\mspage{540}

given by :

\begin{equation*}
\tag{40}
\begin{gathered}
\diagram{\begin{tikzcd}[ampersand replacement=\&]
(x,\underline u)\arrow[r,mapsto]\&(\alpha,\beta,f)
\end{tikzcd}}
\quad\text{with}\quad
\alpha=\underline u\underline a^{-1},\quad
\beta=\underline u\underline b^{-1},\quad
f=\underline u\underline U^{-1}\\
\hfill\text{(how beautiful it is !)}\\
\text{where}\quad\underline a=\unsure{\Theta_\rho}(x),\quad
\underline b=\Theta_\sigma(x),\quad\underline U=\Theta_G(x)
\end{gathered}
\end{equation*}

\deleted{\unsure{The} transformation}\inserted{The operation of $L_0$}
on $\Sigma_{\rho,\sigma}$ by

\begin{equation*}
\tag{41}
\diagram{\begin{tikzcd}[ampersand replacement=\&]
(\alpha,\beta,f)\arrow[r]\&\unsure{''|(}\ell\alpha,\ell\beta,\ell f
\end{tikzcd}}
\end{equation*}

corresponds to the operation of \inserted{left} translation by the
subgroup $L_0$ of $S\mathcal M_{\rho,\sigma}$ (cf. (35)), consisting
of the $(1,\ell)$ with $\ell$ in $L_0$, so we recover
$\mathcal M_{\rho,\sigma}\simeq S\mathcal M_{\rho,\sigma}/L_0$
as the quotient \inserted{(group)} scheme of
$S\Sigma_{\rho,\sigma}^{*}$ by the action (41) of $L_0$.

\uline{Remark}\quad In this description of the $(\alpha,\beta,f)$
the new quantities $\underline u$, $\underline a$,
$\underline b$, $\underline U$ occur, which leads to considering the morphism

\begin{equation*}
\tag{42}
\left\{
\diagram{\begin{tikzcd}[ampersand replacement=\&]
S\mathcal M_{\rho,\sigma}\arrow[r]\&B'_0\times T_\rho\times\unsure{N'_\sigma}\times G\\
(x,\underline u)\arrow[r,mapsto]\&(\underline u,\underline a,\underline b,\underline U)
\end{tikzcd}}\right.
\quad\text{with}\quad
\begin{gathered}
\underline a=\Theta_\rho(x)\\
\underline b=\Theta_\sigma(x)\\
\underline U=\Theta_G(x)
\end{gathered}
\end{equation*}

We have in this connection :

\uline{Proposition}\quad The homomorphism (42) is a monomorphism, its image consists of
the systems satisfying

\begin{equation*}
\tag{43}
\det\underline u=\det\underline b=\det\underline U\quad\text{i.e. }\mu,
\quad\text{and}\quad\det\underline a/\mu\in\Gamma\mu_2
\end{equation*}

This is an immediate consequence of the corollary to the proposition on page 536', and of the
definition of $S\mathcal M_{\rho,\sigma}$ as a fiber product.

\mspage{474}

\scanpage{476}

10\textsuperscript{o}) The section-subgroups
$\mathcal M_{\rho,\sigma}(\sigma)$, $\mathcal M_{\rho,\sigma}(\rho)$.

We define $\mathcal M_{\rho,\sigma}(\sigma)$ by the condition

\[
\beta_0=1\quad\text{i.e.}\quad\beta=f\quad\text{i.e. cf. }\unsure{(40)})
\]

\begin{equation*}
\tag{43}
\underline U=\underline b
\end{equation*}

It is evident from the proposition on page 536' that the subgroup

\begin{equation*}
\tag{44}
\diagram{\begin{tikzcd}[ampersand replacement=\&]
\mathcal M_{\rho,\sigma}(\sigma)\arrow[d,phantom,"\Vert\,\unsure{\mathrm{df}}"]
\subset\mathcal M_{\rho,\sigma})\\
\bigl\{(\underline a,\underline b,\underline U)\in\Gamma\mathcal M_{\rho,\sigma}\mid U=\underline b\bigr\}
\end{tikzcd}}
\end{equation*}

provides a section -- when $\underline a,\underline b$ is given
\inserted{in $\mathcal T_{\rho,\sigma}$},
\inserted{i.e.} satisfying the condition

\[
\unsure{\delta_\rho}(\underline a)/\delta_\sigma(\underline b)\in\Gamma\mu_2,
\]

and the simplest way of lifting it to an element
\inserted{$(\underline a,\underline b,\underline U)$}
of $\mathcal M_{\rho,\sigma}$ isn't it by
$(\underline a,\underline b,\underline U\,\unsure{\mathrm{df}}=\underline b)$ !
(And we \unsure{will have} indeed $\delta_G(\underline U)=\delta_\sigma(\underline b)$ .).
We likewise define $\mathcal M_{\rho,\sigma}(\rho)$
by $\alpha_0=1$ i.e. $\alpha=f$ i.e.

\begin{equation*}
\tag{45}
\diagram{\begin{tikzcd}[ampersand replacement=\&]
\mathcal M_{\rho,\sigma}(\rho)\arrow[d,phantom,"\Vert"]\subset\mathcal M_{\rho,\sigma}\\
\bigl\{(\underline a,\underline b,\underline U)\text{ \unsure{in} }\mathcal M_{\rho,\sigma}\mid\underline U=\underline a\bigr\}
\end{tikzcd}}
\end{equation*}

which gives a partial section \inserted{over} the subgroup
$\mathcal T'_{\rho,\sigma}$ of $\mathcal T_{\rho,\sigma}$ defined as

\begin{equation*}
\tag{46}
\diagram{\begin{tikzcd}[ampersand replacement=\&]
\mathcal T'_{\rho,\sigma}\arrow[d,phantom,"\Vert\,\unsure{\mathrm{df}}"]\unsure{\in|\subset}\mathcal T_{\rho,\sigma}\\
\bigl((\underline a,\underline b)\in T_\rho\times N'_\sigma\mid\unsure{\delta_\rho}(\underline a)=\delta_{\unsure{\sigma|\rho}}(\underline b)\bigr)
\arrow[d,phantom,"\Vert"]\\
\operatorname{Ker}\bigl(\mathcal T_{\rho,\sigma}\xrightarrow{\epsilon_3}\mu_2\bigr)
\end{tikzcd}}
\end{equation*}

so

\begin{equation*}
\tag{47}
\mathcal T'_{\rho,\sigma}=T_\rho\times_{\mathbb{G}_m}\unsure{N'_\sigma},\quad\text{and we have}
\end{equation*}

\mspage{475}

\scanpage{477}
\mspage{541}

so

\begin{equation*}
\tag{48}
\diagram{\begin{tikzcd}[ampersand replacement=\&]
\mathcal T'_{\rho,\sigma}/\mathcal T^0_{\rho,\sigma}\arrow[r,"\epsilon_2"]\&\mu_2
\end{tikzcd}}\,.
\end{equation*}

\uline{Remark}\quad These subgroups can be defined, and provide a
section resp. a section \deleted{on}\inserted{over}
$\mathcal T'_{\rho,\sigma}$, as soon as we have (in the general situation
with $\unsure{7^{\circ}|\mathfrak f}$, page 526' ff) embeddings
\diagram{\begin{tikzcd}[ampersand replacement=\&]
T_\rho\arrow[r,hook]\&G
\end{tikzcd}},
\diagram{\begin{tikzcd}[ampersand replacement=\&]
N'_\sigma\arrow[r,hook]\&G
\end{tikzcd}}, such that
$\delta_\rho,\delta_\sigma$ are induced by $\delta_G$.

In the modular case, we also write

\begin{equation*}
\tag{49}
\mathcal M_{\rho,\sigma}(1/2)=\mathcal M_{\rho,\sigma}(\sigma),\quad
\mathcal M_{\rho,\sigma}(-\bar j)=\mathcal M_{\rho,\sigma}(\rho),
\end{equation*}

and we define more \unsure{general\textsuperscript{ly}|generally}

\begin{equation*}
\tag{50}
\mathcal M_{\rho,\sigma}(q),\qquad
q\in\{\unsure{0},1,\infty\};1/2,\unsure{\underline 2|2},-1\}
\end{equation*}

starting from $\mathcal M_{\rho,\sigma}(0)$ and
$\mathcal M_{\rho,\sigma}(1/2)$ by

\begin{equation*}
\tag{51}
\operatorname{int}(\rho_G)\mathcal M_{\rho,\sigma}(q)
=\mathcal M_{\rho,\sigma}(\dot\rho(q))
\end{equation*}

\margin{$\rho_G$ is $\rho$ as an element of the subgroup
$\unsure{\mathcal G}_{\rho,\sigma}^{+}\simeq SG$ of
$\underset{\unsure{\sim}}{\mathcal M}_{\rho,\sigma}$}

$\dot\rho\deleted{(q)}$ is the homographic transformation

\begin{equation*}
\tag{52}
\dot\rho:\quad
\diagram{\begin{tikzcd}[ampersand replacement=\&]
q\arrow[r,mapsto]\&\dfrac{1}{1-q}
\end{tikzcd}}
\end{equation*}

(whose third iterate is the identity, just like the third iterate
of $\widetilde\rho$).

I would like : to define $\mathcal M_{\rho,\sigma}(-j)$, and
\deleted{\unsure{as}} we have a choice within a well-determined
$\underline{\unsure{\mathrm{Sl}|\mathrm{SL}}(2,\mathbb Z)}^{\wedge}$-conjugacy class,
a choice more

\mspage{476.}

\scanpage{478}
\unsure{restricted} would be between
\[
\operatorname{int}(\sigma_i)(\mathcal M_{\rho,\sigma}(0))
\qquad i\in\{0,1,\infty\}
\tag{53}
\]
\mspage{477}

\scanpage{479}
\mspage{542}
\uline{$\overline{\mathrm{XII}}$}\qquad\uline{Retrospective}

After many unnecessarily complicated detours, I ended up quite simply constructing two group schemes over $\mathbb Z$
\[
\mathcal M_{\rho,\sigma}\subset
\deleted{\underset{\substack{\Vert\\\unsure{\mathrm{Gl}(2)}}}{G}\times}
T_\rho\times N'_\sigma\times G
\tag{1}
\]
\[
S\mathcal M_{\rho,\sigma}\subset
\unsure{B'_0}\times T_\rho\times N'_\sigma\times
\unsure{\deleted{B_0}G}
\tag{2}
\]
by
\[
\mathcal M_{\rho,\sigma}=
\left\{\overbrace{\left(
\unsure{\deleted{\underline U}\underline a},
\unsure{\deleted{\underline a}\underline b},
\unsure{\deleted{\underline b}\underline U}
\right)}^{x}\,\middle|\,
\overbrace{\det\underline U=\det\underline b}^{\chi(x)}=
\overset{\substack{\varepsilon_3(x)\\\vert}}{\varepsilon'}\det\underline a,
{\varepsilon'}^2=1\right\}
\tag{3}
\]
\[
S\mathcal M_{\rho,\sigma}=
\left\{(\unsure{\underline u},\underline a,\underline b,\underline U)\,\middle|\,
\det\underline U=\det\underline u=\det\underline b=
\varepsilon'\det\underline a,\ {\varepsilon'}^2=1\right\}
\tag{4}
\]
\margin{thus $S\mathcal M_{\rho,\sigma}=B'_0\times_{\mathbb{G}_m}\mathcal M_{\rho,\sigma}$}

with
\[
\mathcal M_{\rho,\sigma}\simeq S\mathcal M_{\rho,\sigma}/L_0
\tag{5}
\]
where
\diagram{\begin{tikzcd}[ampersand replacement=\&,column sep=large]
L_0\arrow[r,hook]\&S\mathcal M_{\rho,\sigma}\\
\lambda\arrow[r,mapsto]\&(1,1,1,\lambda)
\end{tikzcd}}

and
\[
S\mathcal M_{\rho,\sigma}\simeq
\underset{\text{}\frac12\text{-direct product}}{\mathcal M_{\rho,\sigma}\mathbin{.}L_0}
\tag{6}
\]
where
\[
\tag{7}
\left\{\vcenter{\hbox{\diagram{\begin{tikzcd}[ampersand replacement=\&,column sep=large]
\mathcal M_{\rho,\sigma}\arrow[r,hook]\&S\mathcal M_{\rho,\sigma}\\
x=(\underline U,\underline a,\underline b)\arrow[r,mapsto]\&
(\underline U,\underline a,\underline b,
 i_{T_0}(\unsure{\deleted{\det}\mu})\deleted{\unsure{(\underline U)}})
\end{tikzcd}}}}\right.
\qquad\text{where }\mu=\chi(x).
\]

Here (for reference) in the modular case
\[
\rho=\begin{pmatrix}0&1\\-1&1\end{pmatrix}
\qquad
\sigma=\begin{pmatrix}0&1\\-1&0\end{pmatrix}
\tag{8}
\]
\[
\left\{\begin{aligned}
T_\rho&=\left\{(c\rho+d)\,\middle|\,
\underbrace{\det(c\rho+d)}_{c^2+cd+d^2}\ \mathrm{inv.}\right\}\\
T'_\sigma&=\left\{(t\sigma+u)\,\middle|\,
\underbrace{\det(t\sigma+u)}_{t^2+u^2}\ \mathrm{inv.}\right\}.
\end{aligned}\right.
\tag{9}
\]
\margin{$\left\{\begin{aligned}
B'_0&=\left\{\begin{pmatrix}\mu&\lambda\\0&1\end{pmatrix}\,\middle|\,\mu\ \mathrm{inv.}\right\}\\
T_0&=\left\{\begin{pmatrix}\mu&0\\0&1\end{pmatrix}\,\middle|\,\mu\ \mathrm{inv.}\right\}\\
i_{T_0}(\mu)&=\begin{pmatrix}\mu&0\\0&1\end{pmatrix}\\
L_0&=\left\{\begin{pmatrix}1&\lambda\\0&1\end{pmatrix}\right\}
\end{aligned}\right.$}
\deleted{(\unsure{10})}
\[
N'_\sigma=\left\{\underline b\text{ in }G\,\middle|\,
\underline b(\sigma)=\underset{\substack{\Vert\\\unsure{\varepsilon_\sigma}(\underline b)}}{\varepsilon}\sigma,
\ \varepsilon^2=1\right\},
\tag{10}
\]
whence exact sequence
\[
\tag{11}
\vcenter{\hbox{\diagram{\begin{tikzcd}[ampersand replacement=\&,column sep=large]
1\arrow[r,dash]\&T_\sigma\arrow[r]\&N'_\sigma\arrow[r,"{\unsure{\varepsilon_\sigma}}"]\&\mu_2\arrow[r]\&1
\end{tikzcd}}}},
\]
\mspage{478}

\scanpage{480}
and the three homomorphisms
\[
\tag{12}
\vcenter{\hbox{\diagram{\begin{tikzcd}[ampersand replacement=\&,column sep=large]
\mathcal M_{\rho,\sigma}\arrow[r,"\Theta_G"]\&G
\end{tikzcd}}}},\quad
\vcenter{\hbox{\diagram{\begin{tikzcd}[ampersand replacement=\&,column sep=large]
\mathcal M_{\rho,\sigma}\arrow[r,"\Theta_\rho"]\&T_\rho
\end{tikzcd}}}},\quad
\vcenter{\hbox{\diagram{\begin{tikzcd}[ampersand replacement=\&,column sep=large]
\mathcal M_{\rho,\sigma}\arrow[r,"\Theta_\sigma"]\&N'_\sigma
\end{tikzcd}}}}
\]
given in the definition (1)(3) are completed by
\[
\tag{12 bis}
\vcenter{\hbox{\diagram{\begin{tikzcd}[ampersand replacement=\&,column sep=large]
\mathcal M_{\rho,\sigma}\arrow[r,"{\chi=\det\circ\Theta_G}"{yshift=3pt}]\&\mathbb{G}_m
\end{tikzcd}}}},\quad
\vcenter{\hbox{\diagram{\begin{tikzcd}[ampersand replacement=\&,column sep=large]
\mathcal M_{\rho,\sigma}\arrow[r,"{\varepsilon_2=\varepsilon_\sigma\circ\Theta_\sigma}"{yshift=3pt}]\&\mu_2
\end{tikzcd}}}},\quad
\vcenter{\hbox{\diagram{\begin{tikzcd}[ampersand replacement=\&,column sep=large]
\mathcal M_{\rho,\sigma}\arrow[r,"{\varepsilon_3(x)=\frac{\det\underline a}{\det\underline b}}"{yshift=3pt}]\&\mu_2
\end{tikzcd}}}}
\]
so that we find
\[
\tag{13}
\vcenter{\hbox{\diagram{\begin{tikzcd}[ampersand replacement=\&,column sep=large]
\mathcal M_{\rho,\sigma}\arrow[r,hook,"{(\Theta_G,\Theta_\rho,\Theta_\sigma,\chi,\varepsilon_2,\varepsilon_3)}"{yshift=5pt}]\&
G\times T_\rho\times N'_\sigma\times\mathbb{G}_m\times\mu_2\times\mu_2\\
x\arrow[r,mapsto]\&(\underline U,\underline a,\underline b,\mu,\varepsilon,\varepsilon')
\end{tikzcd}}}}
\]
whose image consists of the systems \unsure{satisfying}
\[
\left\{\begin{aligned}
&\det\underline U=\det\underline b=\varepsilon'\det\underline a=\mu\\
&\varepsilon=\varepsilon_\sigma(\underline b)\quad\text{i.e}\quad
\begin{gathered}
\underline b(\sigma)=\varepsilon\sigma\quad\text{i.e.}\\
\underbrace{\sigma(\underline b)}_{{}^tb^{-1}}=\varepsilon\underline b
\quad\text{i.e}\quad \underline b\,{}^tb=\varepsilon
\end{gathered}
\end{aligned}\right.
\tag{14}
\]

Likewise we have a monomorphism
\[
\tag{15}
\vcenter{\hbox{\diagram{\begin{tikzcd}[ampersand replacement=\&,column sep=large]
S\mathcal M_{\rho,\sigma}\arrow[r,hook]\&
\unsure{B'_0}\times T_\sigma\times N'_\sigma\times
\underline{\unsure{\deleted{B_0}G}}\times\mathbb{G}_m\times\mu_2\times\mu_2\\
x\arrow[r,mapsto]\&(\unsure{\underline u},\underline a,\underline b,\underline U,\mu,\varepsilon,\varepsilon')
\end{tikzcd}}}}
\]
whose image consists of the systems \unsure{satisfying}
\[
\left\{\begin{aligned}
\mu&=\det\underline u=\det\underline U=\det\underline b=\varepsilon'\det\underline a\\
\varepsilon&=\varepsilon_\sigma(\underline b)\quad\text{i.e}\qquad
\underline b(\sigma)=\varepsilon(\sigma)
\end{aligned}\right.
\tag{16}
\]

Let \deleted{\illegible{1}} on the other hand
\[
\tag{17}
\begin{gathered}
\vcenter{\hbox{\diagram{\begin{tikzcd}[ampersand replacement=\&,column sep=large]
S\Sigma^*_{\rho,\sigma}\arrow[r,hook]\&
G\times\deleted{\unsure{G}}G\times G\times\mathbb G_a\times\mu_2
\end{tikzcd}}}}\\[-2pt]
\Vert\\[-2pt]
\left\{\underbrace{(\alpha,\beta,f,\mu,\varepsilon)}_y\,\middle|\,
\begin{gathered}
(\det\alpha)^2=\det\beta=\det f=1\\
\underbrace{\alpha\rho(\alpha)^{-1}}_\xi=
\underbrace{\varepsilon\beta\sigma(\beta)^{-1}}_{\eta}
\varepsilon_1(\mu-1)
\end{gathered}\right\}
\end{gathered}
\]
(where we define
\[
\varepsilon_1=\begin{pmatrix}1&0\\1&1\end{pmatrix},\qquad
\varepsilon_1(\nu)=\varepsilon_1^\nu\overset{\unsure{\mathrm{dfn}}}{=}
\begin{pmatrix}1&0\\\nu&1\end{pmatrix}.)
\tag{18}
\]
\mspage{479}

\scanpage{481}
\mspage{543}
in fact the projection
\[
\tag{19}
\left\{\vcenter{\hbox{\diagram{\begin{tikzcd}[ampersand replacement=\&,column sep=large]
S\Sigma^*_{\rho,\sigma}\arrow[r,hook]\&G\times G\times G\\
y\arrow[r,mapsto]\&(\alpha,\beta,f)
\end{tikzcd}}}}\right.
\]
is already a monomorphism. We let act \deleted{$L_0$}\inserted{freely}
\[
L_0=\operatorname{Ker}\left(
\vcenter{\hbox{\diagram{\begin{tikzcd}[ampersand replacement=\&,column sep=large]
B'_0\arrow[r,"\det"]\&\mathbb{G}_m
\end{tikzcd}}}}\right)
=\deleted{\unsure{\{}}\left\{\begin{pmatrix}1&\lambda\\0&1\end{pmatrix}\right\}
\quad\text{on }S\Sigma^*_{\rho,\sigma}
\]
by
\[
\tag{20}
\vcenter{\hbox{\diagram{\begin{tikzcd}[ampersand replacement=\&,column sep=large]
(\alpha,\beta,f)\arrow[r,mapsto]\&(l\alpha,l\beta,lf)
\end{tikzcd}}}},
\]
and we set
\[
\Sigma^*_{\rho,\sigma}\overset{\unsure{\mathrm{dfn}}}{=}
\deleted{S\Sigma^*_{\rho,\sigma}}\ L_0\backslash S\Sigma^*_{\rho,\sigma}.
\tag{21}
\]

This being so, we have an \uline{isomorphism}
\[
\tag{22}
\vcenter{\hbox{\diagram{\begin{tikzcd}[ampersand replacement=\&,column sep=large]
S\mathcal M_{\rho,\sigma}\arrow[r,"\sim"]\&S\Sigma^*_{\rho,\sigma}\\
(\underline u,\underline a,\underline b,\underline U)\arrow[r,mapsto]\&
\left\{\begin{gathered}
(\alpha,\beta,f),\\
\alpha=\underline u\,\underline a^{-1},\quad
\beta=\underline u\,\underline b^{-1},\quad
f=\underline u\,\underline U^{-1}
\end{gathered}\right.
\end{tikzcd}}}}
\]
transforming the action of $L_0$ on $S\unsure{\deleted{\Sigma^*}\mathcal M}_{\rho,\sigma}$ \inserted{by \unsure{translation} on the left} \inserted{into} the action of $L_0$ described in (20), and inducing an isomorphism
\[
\tag{23}
\vcenter{\hbox{\diagram{\begin{tikzcd}[ampersand replacement=\&,column sep=large]
\mathcal M_{\rho,\sigma}\arrow[r,"\sim"]\&\Sigma^*_{\rho,\sigma}
\end{tikzcd}}}}
\]
by passage to the quotient.

We have canonical homomorphisms
\[
\tag{24}
\vcenter{\hbox{\diagram{\begin{tikzcd}[ampersand replacement=\&,column sep=large]
S_0\mathcal M_{0,3}^{\sim}\arrow[r]\&S\mathcal M_{\rho,\sigma}(\widehat{\mathbb Z})/\Sigma
\end{tikzcd}}}}
\]
\[
\tag{25}
\vcenter{\hbox{\diagram{\begin{tikzcd}[ampersand replacement=\&,column sep=large]
\mathcal M_{0,3}^{\sim}\arrow[r]\&\mathcal M_{\rho,\sigma}(\widehat{\mathbb Z})/\Sigma
\end{tikzcd}}}}
\]
the second deduced from the first by passage to the quotient, and coinciding on $\mathcal G_{0,3}^{+\wedge}$ with
\deleted{\unsure{and|the} \illegible{2} \unsure{in groups}}
\mspage{480}

\scanpage{482}
the composite homomorphism
\[
\tag{26}
\vcenter{\hbox{\diagram{\begin{tikzcd}[ampersand replacement=\&,column sep=large]
\mathcal G_{0,3}^{+\wedge}
\arrow[r,dash,"{\substack{\S42\ \overline{\underline{\mathrm{IV}}}\\\sim}}"]\&
\underset{\substack{\Vert\\SG(\widehat{\mathbb Z})/SG(\widehat{\mathbb Z})\cap\Sigma}}
{\mathrm{Sl}(2,\widehat{\mathbb Z})'}
\arrow[r,hook]\&
\deleted{\unsure{\mathrm{Sl}(2,\widehat{\mathbb Z})'}}
\mathcal M_{\rho,\sigma}(\widehat{\mathbb Z})/\Sigma
\end{tikzcd}}}}
\]
thanks to the canonical homomorphism
\[
\tag{27}
\vcenter{\hbox{\diagram{\begin{tikzcd}[ampersand replacement=\&,column sep=large,row sep=small]
SG\arrow[r,hook]\arrow[d,equal]\&\mathcal M_{\rho,\sigma}\\
\mathrm{Sl}(2)\deleted{\unsure{\mathbb Z}}\arrow[r,equal]\&
\left\{(\unsure{\underline a},\underline b,\underline U)
\,\middle|\,\overset{\inserted{\text{in }\mathcal M_{\rho,\sigma}}}
{\underline a=\underline b=1}\right\}
\end{tikzcd}}}}
\]

We have defined\deleted{\unsure{t}}
\[
\tag{28}
\vcenter{\hbox{\diagram{\begin{tikzcd}[ampersand replacement=\&,column sep=large]
\Sigma=\underset{\substack{\Vert\\(\mu_2)^3(\mathbb Z)}}{(\mathbb Z/2\mathbb Z)^3}
\arrow[r,hook]\&\mathcal M_{\rho,\sigma}(\mathbb Z)
\arrow[r,phantom,"\subset"]\&\mathcal M_{\rho,\sigma}(\widehat{\mathbb Z})
\end{tikzcd}}}}
\]
via the canonical homomorphism (central)
\[
\tag{29}
\vcenter{\hbox{\diagram{\begin{tikzcd}[ampersand replacement=\&,column sep=large]
\mu_2^3\arrow[r]\&\deleted{\unsure{\mathcal T}}\mathcal M_{\rho,\sigma}
\arrow[r,phantom,"\subset"]\&T_\rho\times N'_\sigma\times G
\end{tikzcd}}}}
\]
defined via the canonical composites
\[
\tag{30}
\vcenter{\hbox{\diagram{\begin{tikzcd}[ampersand replacement=\&,column sep=large,row sep=small]
\&\&T_\rho\&\\
\mu_2\arrow[r,hook]\&\mathbb{G}_m\arrow[ur,hook]\arrow[r,hook]\arrow[dr,hook]\&
T_\sigma\arrow[r,hook]\&N'_\sigma\\
\&\&G\&
\end{tikzcd}}}}.
\]

By passage to the quotient, the homomorphism (25) induces
\[
\tag{\unsure{31|34}}
\vcenter{\hbox{\diagram{\begin{tikzcd}[ampersand replacement=\&,column sep=large]
\Gamma_{\deleted{\unsure{0}}0,3}^{\sim}\arrow[r]\&
\mathcal T_{\rho,\sigma}(\widehat{\mathbb Z})/
\underbrace{\{\pm1\}\times\{\pm1\}}_{\Sigma_0}
\end{tikzcd}}}},
\]
where
\[
\tag{32}
\begin{gathered}
\vcenter{\hbox{\diagram{\begin{tikzcd}[ampersand replacement=\&,column sep=large]
\mathcal T_{\rho,\sigma}\arrow[r,hook]\&T_\rho\times N'_\sigma
\end{tikzcd}}}}\\
\Vert\\
\left\{(\underline a,\underline b)\,\middle|\,
\det\underline a/\det\underline b\in\Gamma\mu_2\right\}
\end{gathered}
\]

The giving --- rather the \uline{construction} made by us --- of a \deleted{(24)} group homomorphism (24)
\mspage{481}

\scanpage{483}
\mspage{544}

amounts to the construction of the four group homomorphisms
\[
\diagram{\begin{tikzcd}[ampersand replacement=\&,column sep=huge]
S_0\mathcal M_{0,3}^{\sim}
  \arrow[r, "{\substack{\text{giving }\underline u\\\Theta_{B'_0}^{\pi}}}"]
  \& B'_0(\widehat{\mathbb Z})=\widehat{\mathbb Z}^{*}\mathbin{.}\widehat{\mathbb Z}
\end{tikzcd}}
\qquad\text{(}\tfrac12\text{-direct product)}
\tag{33}
\]
\[
\diagram{\begin{tikzcd}[ampersand replacement=\&,column sep=huge]
\mathcal M_{0,3}^{\sim}
  \arrow[r, "{\substack{\text{giving }\underline U\text{ (mod sign)}\\\Theta_G^{\pi}}}"]
  \&G(\widehat{\mathbb Z})/\pm1=\mathrm{Gl}(2,\widehat{\mathbb Z})/\pm1
\end{tikzcd}}
\tag{34}
\]
\[
\diagram{\begin{tikzcd}[ampersand replacement=\&,column sep=huge]
\Gamma_{0,3}^{\sim}
  \arrow[r, "{\substack{\text{giving }\underline a\text{ (mod sign)}\\\Theta_{\rho}^{\pi}}}"]
  \&T_{\rho}(\widehat{\mathbb Z})/\pm1
\end{tikzcd}}
\quad\begin{gathered}
\text{(adelic units}\\
\text{(mod }\pm1\text{) of the}\\
\text{field }\mathbb Q(\deleted{\unsure{\sqrt[3]{1}}}\inserted{j})\text{)}
\end{gathered}
\tag{35}
\]
\[
\diagram{\begin{tikzcd}[ampersand replacement=\&,column sep=huge]
\Gamma_{0,3}^{\sim}
  \arrow[r, "{\substack{\text{giving }\underline b\text{ (mod sign)}\\\Theta_{\sigma}^{\pi}}}"]
  \&T_{\sigma}(\widehat{\mathbb Z})/\pm1
\end{tikzcd}}
\quad\begin{gathered}
\text{(adelic units}\\
\text{(mod }\pm1\text{) of the}\\
\text{field }\mathbb Q(i)\text{)}
\end{gathered}
\tag{36}
\]
\unsure{linked|unlinked} by the relations (4), which we are going to identify.
We use the subgroup
\[
\begin{aligned}
\mathcal M_{0,3}^{!\sim}[0]^{\unsure{-}}
  &=\text{subgroup of }\mathcal M_{0,3}^{!\sim}\text{ which}\\
  &\quad\text{(leaves invariant the classes of the loop \unsure{subgroups}}\\
  &\quad\text{of }\Pi_{0,3}^{\wedge}\text{, and which)}\\
  &\quad\text{normalizes }L_{0\pi}\text{ (closed subgroup}\\
  &\quad\text{generated by }l_0\text{)}
\end{aligned}
\tag{37}
\]
and we note that
\[
\diagram{\begin{tikzcd}[ampersand replacement=\&]
1\arrow[r]\&L_{0\pi}\arrow[r]\&\mathcal M_{0,3}^{!\sim}[0]\arrow[r]\&\Gamma_{0,3}^{\sim}\arrow[r]\&1
\end{tikzcd}}
\tag{38}
\]
and that
\[
\left\{\begin{aligned}
\mathcal M_{0,3}^{\sim}
  &=\mathcal M_{0,3}^{!\sim}[0]\mathbin{.}
  \underbrace{\mathcal G_{0,3}^{+\wedge}}_{\substack{\text{invariant subgroup}\\\inserted{\text{of }\mathcal M_{0,3}^{\sim}}}}\\
\text{with}\qquad
\mathcal M_{0,3}^{!\sim}[0]\cap\mathcal G_{0,3}^{+\wedge}
  &=L_{0\pi}
\end{aligned}\right.
\tag{39}
\]
--- of course (38) is a consequence of (39).

finally
\[
S_0\mathcal M_{0,3}^{\sim}\simeq
\underbrace{\mathcal M_{0,3}^{!\sim}[0]\mathbin{.}\mathcal G_{0,3}^{+\wedge}}_{\text{(}\tfrac12\text{-direct product)}}
\tag{40}
\]
1\textsuperscript{o}) The homomorphism (33) \inserted{from $S_0\mathcal M_{0,3}^{\sim}$}
\deleted{\unsure{it}} is trivial on the factor $\mathcal G_{0,3}^{+\wedge}$

\mspage{482}

\scanpage{484}

\margin{so the \deleted{definition}\inserted{description}
of (33) amounts to the description of the \unsure{sole homomorphism|first homomorphism}
\diagram{\begin{tikzcd}[ampersand replacement=\&]
\mathcal M_{0,3}^{\sim}[0]\arrow[r]\&B'_0(\widehat{\mathbb Z})
\end{tikzcd}}}
Its \deleted{\unsure{not}} \unsure{\uline{value}} depends only on the pair
$(\alpha,\beta)$ \deleted{\unsure{part|taken}} of the triple $(\alpha,\beta,f)$
which describes an element of $\mathcal M_{0,3}^{\sim}$), and
the \deleted{\unsure{posi}} definition of $\underline u$ in terms of
$\alpha,\beta$ --- and even only in terms
of the image of $\xi=\alpha\rho(\alpha)^{-1}$ in
$\mathrm{Sl}(2,\widehat{\mathbb Z})$
\deleted{\unsure{modulo} $\unsure{\pm1}$, \unsure{or}}
\inserted{and even} only in
$\mathrm{Sl}(2,\widehat{\mathbb Z})'$ ! ---
\margin{For one \unsure{normalizes} by the condition
$\xi\rho(\xi)\rho^2(\xi)=\unsure{\pm1|+1}$ !}
\unsure{to|has} been \unsure{made|made} essentially in \S45
or at least in \S46 $\overline{\underline{\mathrm{III}}}$, --- we have
\[
\underline u:=\begin{pmatrix}1/D&B/D\\0&1\end{pmatrix}
\qquad\begin{gathered}
\text{where }\xi=\begin{pmatrix}A&B\\C&D\end{pmatrix}\\
\text{(in }\mathrm{Sl}(2,\widehat{\mathbb Z})\text{)}
\end{gathered}
\tag{41}
\]
or again, in terms of
\[
\underline\alpha=\begin{pmatrix}a&b\\c&d\end{pmatrix}
\qquad\text{with }ad-bc\unsure{=|\in}\inserted{\varepsilon'\in}\{\pm1\}
\]
\[
\left\{\begin{aligned}
\inserted{\varepsilon'}D&=(c^2+cd+d^2)\deleted{\unsure{/\varepsilon}}\\
\inserted{\varepsilon'}B&=ac+bd+bc
\end{aligned}\right.
\qquad\deleted{\text{with }\unsure{\varepsilon'}}
\tag{42}
\]

This does not depend on the foundations not laid for profinite groups ---
the ``calculations in $\xi,\eta$'' are in fact sufficient.
It is thanks to the \unsure{description} of this homomorphism
that one can define
\[
\begin{gathered}
\mathcal M_{0,3}^{\sim}(0)\subset\mathcal M_{0,3}^{\sim}[0]\\
\begin{array}{c}
\Vert\\
\{\underline{\underline u}\mid\underline u\in T_0(\widehat{\mathbb Z})
\quad\text{i.e. }B=0\}\\
\qquad\text{i.e. }ac+bd+bc=0
\end{array}
\end{gathered}
\tag{43}
\]
(normalization condition),

\mspage{483}

\scanpage{485}
\mspage{545}

and obtain canonical decompositions into
$\tfrac12$-\unsure{direct|direct} products
\[
\mathcal M_{0,3}^{\sim}[0]\simeq\mathcal M_{0,3}^{\sim}(0)\mathbin{.}L_{0\pi}
\tag{44}
\]
\[
\mathcal M_{0,3}^{\sim}\simeq\mathcal M_{0,3}^{\sim}(0)\mathbin{.}\mathcal G_{0,3}^{+\wedge}.
\tag{45}
\]

2\textsuperscript{o}) To describe (34), since it is already known on
$\mathcal G_{0,3}^{+\wedge}$, it suffices to describe it on
$\mathcal M_{0,3}^{\sim}[0]$ or even only on
$\mathcal M_{0,3}^{\sim}(0)$ (subject of course to compatibilities\ldots).
And on $\mathcal M_{0,3}^{\sim}[0]$ it coincides with (33) ---
whereas on $\mathcal M_{0,3}^{\sim}(0)$ it is practically trivial
and reduces to the multiplier $\mu$
\[
\Theta_G^{\pi}=\Theta_{B'_0}^{\pi}
\quad\text{on }\mathcal M_{\unsure{\rho|0},3}^{\sim}[0]
\tag{46}
\]
\[
\left\{\begin{aligned}
\Theta_G^{\pi}(\underline{\underline u})
  &=i'_{T_0}(\mu)
  \qquad\begin{gathered}\text{where }\mu\text{ is the}\\
  \text{multiplier of }\underline{\underline u}\end{gathered}\\
\text{\unsure{if|where}}\quad
\underline{\underline u}&\in\mathcal M_{0,3}^{\sim}(0)
\end{aligned}\right.
\tag{47}
\]
(this is even a possible definition of $\mathcal M_{0,3}^{\sim}(0)$).

\uline{Remark} As $\Theta_G^{\pi}$ was constructed
``by sheer strength of \unsure{wrist|wrists}'' by way
of $\Theta_{B'_0}^{\pi}$, i.e. $\underline U$ was constructed from $\underline u$,
there was for a long time a sort of

\mspage{484}

\scanpage{486}

confusion in my mind between $\underline u$ and $\underline U$,
all the more so since most of the time I restricted myself to working
in $\mathcal M_{\unsure{0|\rho},3}^{\sim}[0]$
(after having fumbled about quite a bit with more
\unsure{varied} points of view, before realizing that technically this was
the most elegant way to proceed --- from \S46 on).
This confusion was cleared up for good only yesterday and today,
by dint of putting my feet in my own dishes again and again\ldots

3\textsuperscript{o}) The homomorphisms (35)(36) were obtained
only by working in depth on the parameters $\alpha,\beta$
\unsure{for} $\underline u$ --- here again, it suffices to work
in $\mathcal M_{0,3}^{\sim}[0]$, and to define
\[
\diagram{\begin{tikzcd}[ampersand replacement=\&,row sep=tiny,column sep=huge]
\underline{\underline u}\arrow[r,mapsto]\&\underline a\\
\mathcal M_{0,3}^{\sim}[0]\arrow[r]\&T_{\rho}(\widehat{\mathbb Z})/\pm1\\
{}\arrow[r]\&\unsure{N'_{\sigma}}(\widehat{\mathbb Z})/\pm1\\
\underline{\underline u}\arrow[r,mapsto]\&\underline b
\end{tikzcd}}
\]

\mspage{485}

\scanpage{487}
\mspage{546}

so that it is trivial on $L_0^{\deleted{\unsure{\xi|\wedge}}\pi}$.

The quantities $\underline a,\underline b$ appeared in the
calculations gradually starting with
\S44 (page 366$'$ and following), light began to be shed on their
role only
starting with \S48 $\overline{\underline{\mathrm{III}}},
\overline{\underline{\mathrm{IV}}},\overline{\underline{\mathrm{V}}}$ ---
and a complete understanding of their nature, from the point
of view of the group schemes that enter into the theory,
likewise came about only in these very last
days. This understanding moreover remains for the moment rather
formal, and one should go deeper
in several directions,

a) Make explicit the meaning from the point
of view of the class fields of
$\mathbb Q(j)$ and \inserted{of} $\mathbb Q(\deleted{j}\inserted{i})$.

b) Make explicit the meaning for
the theory of ``\unsure{congruential|congruential}'' maps
(i.e. \unsure{subordinate|subordinate} to \deleted{\unsure{continuous}} representations
of the quotient $\mathrm{Sl}(2,\widehat{\mathbb Z})'$ of
$\mathrm{Sl}(2,\mathbb Z)^{\wedge\prime}$)
of the results obtained,

\mspage{486}

\scanpage{488}

c) Introduce on $\mathcal M_{\rho,\sigma}$ the elements of structure characteristic of arithmetic Galois groups -- and make the connection with the typical constructions of this sort in class field theory.

d) Examine to what extent in this context one can bring in the seemingly draconian compatibility conditions on $\alpha,\beta$ of §47 linked to compatibility with
\[
\diagram{\begin{tikzcd}[ampersand replacement=\&,column sep=large]
\varphi:\widehat\Pi'_{0,3} \arrow[r,"\sim"] \& \mathcal G_{0,3}^{+\wedge}
\end{tikzcd}}
\]
-- I \unsure{foresee.} moreover that the framework of the preceding calculations is too narrow to manage to bring these additional conditions into it.

e) What interests me most for the moment is a systematization of the preceding calculations in a framework much broader than that of the group $SG$ alone and of the \unsure{standard} homomorphism

\mspage{487}

\scanpage{489}
\mspage{547}

\[
\diagram{\begin{tikzcd}[ampersand replacement=\&,column sep=large]
\mathcal G_{0,3}^{+\wedge} \arrow[r] \& \deleted{\unsure{S}}G(\widehat{\mathbb Z})/\pm1
\end{tikzcd}}
\]
-- I have the feeling of having put my finger on a systematic procedure for constructing ``motivic'' representations of $\Gamma_{\mathbb Q}$, via representations of the ``geometric'' part $\mathcal G_{0,3}^{+\wedge}$ of
\[
\mathcal M_{0,3}^{+}\simeq\Pi_1(M_{0,3\,\mathbb Q},\unsure{\mathfrak S}_3)
\simeq\Pi_1(M'_{1,1\,\mathbb Q})
\]
-- or better, \inserted{via} representations of the discrete group $\mathcal G_{0,3}^{+}\simeq\mathrm{Sl}(2,\mathbb Z)'$ (or the \unsure{``}variants such as $\mathrm{Sl}(2,\mathbb Z)$, $\Pi_{0,3}$ \ldots) with values in the integral points of group schemes (\inserted{\unsure{generically}} reductive ??) over $\mathbb Z$.

f) Yet there is another direction of investigation that attracts me a great deal, not necessarily linked to the introduction of group schemes, and which could be pursued \inserted{\unsure{in vitro|in addition}} within the framework of just the ``big'' profinite fundamental groups --

\mspage{488}

\scanpage{490}

namely the study of certain special \inserted{$\mathbb Q$-}coverings of $\unsure{U}_{0,3\,\mathbb Q}=\mathbb P_{\unsure{0,3}}-\{0,1,\infty\}$, and information one should draw from them about additional properties of the ``arithmetic'' outer actions of $\Pi_{0,3}^{\wedge}$ (of the type of the results obtained in §47), and \deleted{\unsure{on}} calculations of splittings of the extension $\mathcal M_{0,3}$ of $\Gamma_{0,3}$ by $\mathcal G_{0,3}^{+\wedge}$ at points other than the ``critical'' points $0,1,\infty,\frac12,2,-1,-\bar j,\unsure{-j}$.

(Not counting the question \unsure{of establishing|of securing} the foundations, via the method suggested by Deligne -- and of studying the loop equation in the \unsure{protilp.|pro. tilp.} compactification $\unsure{L_A|L_\Delta}$ -- and of continuing to investigate in the direction of the \unsure{fundamental} conjectures, notably on the \unsure{arithmetization} of \uline{complex} points -- etc etc).

\mspage{489}

\scanpage{491}
\mspage{548}

Let us return for the moment to the question of identifying the \inserted{four} homomorphisms obtained (33) to (36), in terms of classical things. One can consider the homomorphism (33) \inserted{$\Theta_{B'_0}^{\pi}$} as being of an essentially technical nature, or at least as being subordinate to the homomorphism $\Theta_G^\pi$ of (34). Moreover, (35) and (36) are deduced from it, by composing \inserted{them} with the two homomorphisms
\[
\diagram{\begin{tikzcd}[ampersand replacement=\&,column sep=large]
\Gamma_{0,3}^{\sim} \arrow[r] \& \mathcal M_{0,3}^{\sim}
\end{tikzcd}}
\]
\inserted{(morally at least, for the moment !)} associated with the \deleted{\unsure{the}} points $\frac12$ and $-\bar j$ of $\unsure{U}_{0,3\,\mathbb Q}(\mathbb C)$, and which correspond, at the level of group schemes, to the two splittings $\psi_\rho,\psi_\sigma$ of $\overset{\inserted{(3)}}{\mathcal M_{\rho,\sigma}}$ over $\overset{\inserted{(31)}}{\mathcal T_{\rho,\sigma}}$, given \deleted{by} respectively by
\diagram{\begin{equation*}
\tag{48}
\left\{\begin{gathered}
\begin{tikzcd}[ampersand replacement=\&,row sep=small,column sep=large]
\psi_\rho:(\underline a,\underline b) \arrow[r,mapsto] \& (\underline a,\underline b,\underline a) \\
\psi_\sigma:(\underline a,\underline b) \arrow[r,mapsto] \& (\underline a,\underline b,\underline b)
\end{tikzcd}
\end{gathered}\right.
\end{equation*}}
\margin{\diagram{\begin{tikzcd}[ampersand replacement=\&,row sep=small,column sep=large]
\deleted{\unsure{\mathcal T_{\rho,\sigma}|\psi_\rho}}\,\mathcal T'_{\rho,\sigma}\arrow[r] \& \mathcal M_{\rho,\sigma} \\
\mathcal T_{\rho,\sigma}\arrow[r] \& \mathcal M_{\rho,\sigma},
\end{tikzcd}}}
with this grain of salt that the first is defined only on the subgroup $\mathcal T'_{\rho,\sigma}$ defined by
\[
\tag{49}
\left\{\begin{aligned}
\mathcal T'_{\rho,\sigma}&\subset\mathcal T_{\rho,\sigma}\\
\Vert\quad&\left\{\left(\underset{\inserted{\substack{\rotatebox[origin=c]{-90}{$\unsure{\in}$}\\T_\rho}}}{\underline a},\underset{\inserted{N'_\sigma}}{b}\right)\,\middle|\,\deleted{\text{in }\mathcal T_{\rho,\sigma}}\ \det\underline a=\det\underline b\right\}\\
&\Vert\quad T_\rho\mathbin{\underset{\mathbb{G}_m}{\times}}N'_\sigma.
\end{aligned}\right.
\]

\mspage{490}
